\chapter{Zika Virus Testing}

\section*{What Is This Test?}
Zika virus testing detects Zika virus (ZIKV), a single-stranded positive-sense RNA flavivirus transmitted primarily by \textit{Aedes aegypti} and \textit{A. albopictus} mosquitoes. Additional modes of transmission include sexual (vaginal, anal, oral), vertical (maternal-fetal transplacental), blood transfusion, and laboratory exposure. Most ZIKV infections (60--80\%) are asymptomatic. Symptomatic Zika typically causes a mild, self-limited illness: acute onset of low-grade fever, maculopapular rash (spreading from face to trunk and extremities), non-purulent conjunctivitis (bilateral, itching), arthralgia (particularly small joints of hands and feet), myalgia, and headache. Symptoms last 2--7 days and are often indistinguishable from dengue and chikungunya virus infections (which are transmitted by the same mosquito vectors and co-circulate in many regions). The devastating complication that brought Zika to global attention is congenital Zika syndrome (CZS) resulting from maternal-fetal transmission during pregnancy. CZS features include microcephaly (typically severe, with collapsed skull and overlapping cranial sutures), intracranial calcifications (especially at the gray-white matter junction), ventriculomegaly, cortical thinning, cerebellar hypoplasia, arthrogryposis, ocular abnormalities (macular scarring, chorioretinal atrophy), and sensorineural hearing loss. Zika is also associated with Guillain-Barré syndrome (GBS) in adults. Following the 2015--2016 pandemic in the Americas, Zika transmission has declined dramatically, but focal outbreaks continue. Nucleic acid testing (RT-PCR) of serum and urine in the acute phase plus serologic testing (Zika IgM by MAC-ELISA) is the standard diagnostic approach, complicated by extensive flavivirus cross-reactivity.

\section*{Alternative Names}
\begin{itemize}
  \item Zika PCR
  \item Zika Virus RNA
  \item Zika Serology
  \item Zika IgM
  \item Zika MAC-ELISA
  \item ZIKV Testing
  \item Zika NAT
  \item Flavivirus Serology
\end{itemize}
\section*{Principle}
RT-PCR: Real-time reverse-transcription PCR detecting ZIKV RNA in serum (collected within 7 days of symptom onset), urine (within 14 days), whole blood, CSF, or amniotic fluid. FDA-cleared multiplex assays (e.g., cobas Zika, Aptima Zika, Trioplex real-time RT-PCR targeting ZIKV, DENV, and CHIKV) are used. Sensitivity is highest during the first 7 days of symptoms; thereafter, viremia declines and serology is the primary diagnostic tool. Urine PCR may remain positive for a longer duration (up to 14--21 days after symptom onset) and is recommended as an adjunct to serum PCR. IgM antibody capture ELISA (MAC-ELISA): The CDC Zika MAC-ELISA detects ZIKV IgM antibodies in serum or CSF. IgM typically becomes detectable 4--7 days after symptom onset and persists for weeks to months. However, significant cross-reactivity with other flaviviruses (dengue, yellow fever, West Nile, Japanese encephalitis, St. Louis encephalitis, and tick-borne encephalitis viruses) is a major limitation. A positive ZIKV IgM MAC-ELISA result is "presumptive" and must be confirmed by the plaque reduction neutralization test (PRNT), which measures virus-specific neutralizing antibodies. PRNT: The gold standard for flavivirus serologic confirmation but labor-intensive, slow (1--2 weeks), and still not perfectly specific (secondary flavivirus infections can cross-neutralize due to original antigenic sin). PRNT titers $\geq$1:10 or $\geq$1:20 (depending on the assay and the virus) against ZIKV, with a $\geq$4-fold higher titer than against other flaviviruses (DENV-1--4, WNV, YFV), confirm ZIKV infection. PRNT is essential for distinguishing ZIKV from DENV in regions where both circulate. IgG avidity testing: IgG avidity assays can distinguish primary ZIKV infection (low avidity) from past infection, which is particularly useful for pregnant women with possible remote exposure.

\section*{History}
Zika virus was first isolated from a sentinel rhesus macaque in the Zika Forest of Uganda in 1947. For decades, ZIKV was an obscure pathogen causing sporadic human infections and small outbreaks in Africa and Asia. The first large outbreak occurred on Yap Island, Federated States of Micronesia, in 2007. The association with GBS was noted during the 2013--2014 outbreak in French Polynesia. The devastating 2015--2016 pandemic swept through the Americas (starting in Brazil), with the link to microcephaly recognized in October 2015 and the WHO declaring a Public Health Emergency of International Concern (PHEIC) in February 2016. The CDC developed the Trioplex RT-PCR assay and Zika MAC-ELISA during the emergency. The pandemic subsided by 2017 due to herd immunity and vector control. Since 2018, Zika incidence has plummeted globally, though sporadic cases and small outbreaks continue. The FDA recommended universal Zika testing of all blood donations in 2016; this requirement was lifted in 2021 due to the low ongoing risk. No antiviral therapy exists; treatment is supportive. Vaccines are in development but none is licensed as of 2024.

\section*{Why This Test Is Ordered}
\begin{itemize}
  \item Evaluation of acute Zika virus infection in a symptomatic patient (fever, maculopapular rash, conjunctivitis, arthralgia) with recent travel to or residence in a Zika-endemic area or sexual contact with someone who traveled to such an area. Testing should be performed as soon as possible
  \item Testing of pregnant women with possible Zika exposure (residence in or travel to an area with risk of Zika, or sexual contact with someone who traveled to such an area), regardless of symptoms. Testing algorithm per CDC: If the exposure was within the preceding 2--12 weeks, serum and urine NAAT (PCR) and Zika IgM serology should be performed. Testing asymptomatic pregnant women is recommended per CDC guidance for travel-associated exposure
  \item Evaluation of infants with clinical findings suggestive of congenital Zika syndrome (microcephaly, intracranial calcifications, ventriculomegaly, thin cortex, arthrogryposis, ocular abnormalities) born to mothers with possible Zika exposure during pregnancy. Infant testing: Serum and urine NAAT within the first few days of life; serum Zika IgM; CSF Zika IgM and NAAT if CSF is obtained for other reasons; placental tissue NAAT and histopathology; and head ultrasound/MRI
  \item Evaluation of patients with suspected GBS with recent Zika exposure
  \item Public health surveillance and epidemiologic investigations during suspected Zika outbreaks
  \item Pre-conception counseling: Serologic testing for Zika IgG may be considered in women planning pregnancy with a history of possible exposure, but interpretation is complicated by flavivirus cross-reactivity and waning antibody levels over time
\end{itemize}

\section*{Primary vs Secondary Test}
RT-PCR on serum and urine is the primary diagnostic test during the acute phase (first 7 days of symptoms for serum; first 14 days for urine). Zika IgM MAC-ELISA is the primary test after the first week of illness and for evaluating pregnant women. PRNT is the secondary confirmatory test required for positive or equivocal IgM results due to flavivirus cross-reactivity.

\section*{How This Test Is Performed}
Acute phase (days 1--14 from symptom onset): Serum (SST tube, 5--7 mL) and urine (clean-catch, 10--30 mL) are collected for ZIKV RT-PCR. Whole blood (EDTA) may also be tested. Specimens for PCR should be refrigerated at 2--8\textdegree{}C and transported to the laboratory within 48 hours, or frozen at -70\textdegree{}C. Convalescent phase (beyond 7 days): Serum for Zika IgM MAC-ELISA. CSF: If lumbar puncture is performed for other clinical indications, CSF should be tested for Zika IgM and PCR. Amniotic fluid: For pregnant women with confirmed or suspected Zika, amniocentesis at $\geq$15 weeks gestation for ZIKV RT-PCR may be considered, though sensitivity is low and a negative result does NOT exclude fetal infection. Placental tissue: If the pregnancy results in fetal loss or delivery, placental tissue can be tested by RT-PCR and histopathology.

\section*{What Preparation Is Needed}
No special patient preparation. For PCR testing, the timing of specimen collection relative to symptom onset is critical: serum PCR should be collected within 7 days; urine PCR within 14 days. Specimens collected outside these windows have substantially reduced sensitivity and serology should be the primary test.

\section*{Contraindications and Precautions}
\begin{itemize}
  \item The most significant limitation of Zika diagnostics is the extensive serologic cross-reactivity among flaviviruses (Zika, dengue, West Nile, yellow fever, Japanese encephalitis, St. Louis encephalitis, Powassan). A positive Zika IgM MAC-ELISA, particularly in a patient with prior flavivirus infection or vaccination (dengue, yellow fever vaccine, JE vaccine), cannot reliably distinguish Zika from other flaviviruses. PRNT must be performed for confirmation, and even PRNT may be equivocal in secondary flavivirus infections due to the original antigenic sin phenomenon (anamnestic antibody response directed predominantly against the first-infecting flavivirus). Zika IgM may persist for months, limiting the ability to determine when infection occurred (important for pregnancy management). A positive Zika IgM in a pregnant woman cannot date the trimester of infection. Zika IgG avidity testing, where available, may help distinguish recent from remote infection. False-negative Zika NAAT is common when specimens are collected >7 days (serum) or >14 days (urine) after symptom onset. In the setting of a negative NAAT and a positive or equivocal IgM, determining the timing of infection relies on a combination of epidemiologic history, clinical presentation, and serology. Pre-test counseling is essential for pregnant women, given the implications of a positive test.
\end{itemize}
\section*{Risks}
Venipuncture, urine collection, and CSF collection carry standard procedural risks. The major clinical risk is the emotional and obstetric impact of a positive Zika test during pregnancy, including decisions about pregnancy continuation. Given the complexities of test interpretation, all positive and equivocal results should be discussed in consultation with maternal-fetal medicine specialists and the CDC or state public health department.

\section*{What Happens After the Test}
RT-PCR results available in 4--24 hours. Zika MAC-ELISA results in 2--5 days (performed at CDC, state public health laboratories, or select commercial reference labs). PRNT results in 1--2 weeks. For symptomatic non-pregnant patients with confirmed Zika: Supportive care only (rest, hydration, acetaminophen---NSAIDs should be avoided until dengue is excluded due to risk of hemorrhage). For pregnant women with confirmed Zika: Serial fetal ultrasound (every 3--4 weeks) to monitor for congenital Zika syndrome. Amniocentesis for ZIKV RT-PCR may be offered. Management decisions involve maternal-fetal medicine and infectious disease specialists. For infants with suspected congenital Zika syndrome: Comprehensive evaluation including head ultrasound/MRI, ophthalmologic exam, auditory brainstem response (ABR) testing, and developmental follow-up.

\section*{Understanding Results}

\subsection*{Below Normal}
\begin{itemize}
  \item \textbf{Zika RT-PCR negative on serum and urine:} No ZIKV RNA detected. A negative PCR, particularly from specimens collected outside the optimal window (>7 days serum, >14 days urine), does NOT exclude Zika. A negative PCR combined with a negative Zika IgM (collected $\geq$4 days after symptom onset) makes acute ZIKV infection very unlikely
  \item \textbf{Zika IgM negative:} No ZIKV IgM detected. If the specimen was collected $\geq$4 days after symptom onset (after the window for IgM seroconversion), this largely excludes acute ZIKV infection. However, specimens collected 0--3 days after symptom onset may be falsely negative because IgM has not yet appeared
  \item \textbf{Zika PRNT negative (titer <1:10):} No neutralizing antibodies against ZIKV detected. This essentially excludes ZIKV infection. If ZIKV IgM was positive, the positive IgM is likely due to cross-reactivity with another flavivirus (dengue, WNV) or a false-positive result
\end{itemize}

\subsection*{Above Normal}
\begin{itemize}
  \item \textbf{Zika RT-PCR positive (serum, urine, whole blood, CSF, amniotic fluid):} Confirms ZIKV infection. In pregnancy: maternal NAAT positive indicates acute ZIKV infection, and the fetus is at risk for congenital Zika syndrome. Amniotic fluid NAAT positive confirms fetal infection; however, a negative amniotic fluid NAAT does not exclude fetal infection due to intermittent viral shedding or sampling at a time of low viral load
  \item \textbf{Zika IgM positive (presumptive):} Indicates recent flavivirus infection. Must be confirmed by PRNT. A positive Zika IgM alone is NOT a confirmed diagnosis and should not be communicated to the patient as definitive. Confirmatory PRNT is essential
  \item \textbf{Zika PRNT positive with titer $\geq$1:10 (or $\geq$1:20) against ZIKV and $\geq$4-fold higher than other tested flaviviruses:} Confirms ZIKV infection. This is considered laboratory-confirmed Zika. Without the $\geq$4-fold differential, the result may represent cross-reactivity from another flavivirus (e.g., dengue) and is "equivocal"
  \item \textbf{Zika IgM positive in CSF:} Indicates intrathecal synthesis of ZIKV-specific antibodies. In a neonate or adult with neurologic symptoms, this is highly specific for CNS ZIKV infection
  \item \textbf{Zika IgG avidity low:} Indicates recent (primary) ZIKV infection (within approximately 3--6 months). High avidity indicates past (more remote) or secondary flavivirus infection. Useful as an adjunct when timing of infection in pregnancy is critical
\end{itemize}

\section*{Normal Results}
\textbf{ZIKV RNA not detected} by RT-PCR (serum, urine). \textbf{Zika IgM negative.} \textbf{Zika PRNT negative} (titer <1:10). No evidence of ZIKV infection.

\section*{Follow-up Tests}
\begin{itemize}
  \item \textbf{Zika PRNT (plaque reduction neutralization test):} Mandatory for all positive or equivocal Zika IgM results to confirm the specificity of the IgM. PRNT against ZIKV and dengue virus serotypes 1--4 (and possibly other flaviviruses based on exposure history) is performed. A $\geq$4-fold higher titer to ZIKV versus dengue confirms Zika; equivocal results may occur, particularly in secondary flavivirus infections
  \item \textbf{Dengue PCR and serology (DENV NS1 antigen, DENV IgM/IgG):} When Zika testing is positive or equivocal, dengue co-infection or cross-reaction should be evaluated. In regions co-endemic for ZIKV and DENV, patients should be tested for both viruses simultaneously. Dengue exclusion is particularly important because dengue has higher risk for severe hemorrhagic complications and requires close monitoring
  \item \textbf{Chikungunya serology (CHIKV IgM/IgG and PCR):} Chikungunya co-circulates with Zika and dengue and causes similar acute illness (fever, rash, arthralgia). Arthralgia is typically more severe and prolonged in chikungunya
  \item \textbf{Serial fetal ultrasound:} For pregnant women with confirmed Zika infection, serial ultrasound imaging every 3--4 weeks to detect fetal abnormalities, particularly microcephaly, intracranial calcifications, ventriculomegaly, and growth restriction
  \item \textbf{Infant evaluation for congenital Zika syndrome:} Head ultrasound/MRI, ophthalmologic (retinal) examination, auditory brainstem response (ABR) hearing screen, and Zika NAAT/IgM on infant serum, urine, and CSF (if obtained) within the first days of life. Long-term neurodevelopmental follow-up is essential
\end{itemize}

\section*{References}
\begin{enumerate}
  \item Oduyebo T, Polen KD, Walke HT, et al. Update: interim guidance for health care providers caring for pregnant women with possible Zika virus exposure --- United States. MMWR Morb Mortal Wkly Rep. 2017;66(31):815-822.
  \item CDC. Zika virus: testing guidance. Centers for Disease Control and Prevention; 2019 (updated periodically).
  \item Petersen LR, Jamieson DJ, Powers AM, Honein MA. Zika virus. N Engl J Med. 2016;374(16):1552-1563.
  \item Brasil P, Pereira JP, Moreira ME, et al. Zika virus infection in pregnant women in Rio de Janeiro. N Engl J Med. 2016;375(24):2321-2334.
  \item Sharp TM, Fischer M, Muñoz-Jordán JL, et al. Dengue and Zika virus diagnostic testing for patients with a clinically compatible illness and risk for infection with both viruses. MMWR Recomm Rep. 2019;68(1):1-10.
\end{enumerate}

\clearpage
\section*{Regulatory Status and Authorization Date}
This chapter describes a test category, not a specific commercial product. FDA, CDSCO, EU IVDR, and MHRA approval or registration dates therefore cannot be assigned without the assay manufacturer, product name, instrument, and intended use. For laboratory-developed tests, an individual agency approval date may not exist. See the Preface for the interpretation of regulatory status.

\clearpage
